\subsection{Hyperplanes, Nearest Points, and Separation}

\begin{definition} \label{def:affine-hyperplanes-halfspaces-separation} \hlt{Hyperplanes, Half-Spaces, and Separation}\\
Let $n\in\mathbb{Z}_{>0}$, $\mathbf{N}\in\R^n\setminus\{\mathbf{0}\}$, and $b\in\R$.
\begin{enumerate}[label=(\roman*)]
\setlength{\itemsep}{0pt}
\item The \hlt{affine hyperplane} with normal $\mathbf{N}$ and level $b$ is $\mathcal{H}(\mathbf{N},b)\defeq\{\mathbf{x}\in\R^n:\innerprod{\mathbf{N},\mathbf{x}}=b\}=\frac{b}{\norm{\mathbf{N}}^2}\mathbf{N}+\operatorname{span}(\mathbf{N})^\perp$.
\item The case $b=0$ is the \hlt{linear hyperplane} $\mathcal{H}(\mathbf{N},0)=\operatorname{span}(\mathbf{N})^\perp$.
\item The \hlt{closed half-spaces} are $\mathcal{H}_+(\mathbf{N},b)\defeq\{\mathbf{x}\in\R^n:\innerprod{\mathbf{N},\mathbf{x}}\geq b\}$ and $\mathcal{H}_-(\mathbf{N},b)\defeq\{\mathbf{x}\in\R^n:\innerprod{\mathbf{N},\mathbf{x}}\leq b\}$.\\
Replacing $\geq,\leq$ by $>,<$ defines the corresponding \hlt{open half-spaces}.
\item For $X\subseteq\R^n$ and $\star\in\{<,\leq,>,\geq\}$, the \hlt{set-inequality convention} is $\innerprod{\mathbf{N},X}\star b\Longleftrightarrow(\forall\,\mathbf{x}\in X)\,\innerprod{\mathbf{N},\mathbf{x}}\star b$.
\item For $\sigma\in\{\pm1\}$, the pair $\sigma(\mathbf{N},b)\defeq(\sigma\mathbf{N},\sigma b)$ is an \hlt{orientation} of $\mathcal{H}(\mathbf{N},b)$.
\item Sets $X,Y\subseteq\R^n$ \hlt{strictly separated} by $\mathcal{H}(\mathbf{N},b)$ iff $(\exists\,\sigma\in\{\pm1\})[(\forall\,\mathbf{x}\in X)\,\innerprod{\sigma\mathbf{N},\mathbf{x}}<\sigma b\land(\forall\,\mathbf{y}\in Y)\,\sigma b<\innerprod{\sigma\mathbf{N},\mathbf{y}}]$.
\item Sets $X,Y\subseteq\R^n$ \hlt{strongly separated} by $\mathcal{H}(\mathbf{N},b)$ iff $(\exists\,\sigma\in\{\pm1\})(\exists\,\varepsilon>0)[(\forall\,\mathbf{x}\in X)\,\innerprod{\sigma\mathbf{N},\mathbf{x}}<\sigma b-\varepsilon\land(\forall\,\mathbf{y}\in Y)\,\sigma b+\varepsilon<\innerprod{\sigma\mathbf{N},\mathbf{y}}]$.
\end{enumerate}
\end{definition}

\begin{remark} \label{rem:hyperplane-basic-geometry} \hlt{Hyperplane Geometry}\\
The identity $\mathcal{H}_+(\mathbf{N},b)\cap\mathcal{H}_-(\mathbf{N},b)=\mathcal{H}(\mathbf{N},b)$ holds, and the two open half-spaces together with $\mathcal{H}(\mathbf{N},b)$ partition $\R^n$.\\
Also $\mathbf{N}\in\mathcal{H}(\mathbf{N},\norm{\mathbf{N}}^2)$, and $(\forall\,\mathbf{x}\in\mathcal{H}(\mathbf{N},\norm{\mathbf{N}}^2))\,\norm{\mathbf{N}}^2=\innerprod{\mathbf{N},\mathbf{x}}\leq\norm{\mathbf{N}}\norm{\mathbf{x}}$ by Theorem \ref{thm:cauchy-schwarz-inequality}(i).\\
Thus $\mathbf{N}$ is a closest point of this hyperplane to $\mathbf{0}$.
\end{remark}

\begin{theorem} \label{thm:nearest-point-and-point-separation-convex-set} \hlt{Nearest Points and Separating a Point}\\
Let $C\subseteq\R^n$ be nonempty, closed, and convex.
\begin{enumerate}[label=(\roman*)]
\setlength{\itemsep}{0pt}
\item $(\exists!\,\mathbf{N}_C\in C)(\forall\,\mathbf{x}\in C)\,\norm{\mathbf{N}_C}\leq\norm{\mathbf{x}}$.
\item If $\mathbf{b}\notin C$, put $D\defeq C-\mathbf{b}\defeq\{\mathbf{x}-\mathbf{b}:\mathbf{x}\in C\}$ and let $\mathbf{N}$ be its unique minimum-norm vector.\\
Then $\mathbf{N}\neq\mathbf{0}$ and $\innerprod{\mathbf{N},C}\geq\norm{\mathbf{N}}^2+\innerprod{\mathbf{N},\mathbf{b}}>\innerprod{\mathbf{N},\mathbf{b}}$.
\item Under the hypotheses and notation of (ii), set $\alpha\defeq\innerprod{\mathbf{N},\mathbf{b}}+\frac{1}{2}\norm{\mathbf{N}}^2$. Then $\mathcal{H}(\mathbf{N},\alpha)$ strongly separates $\{\mathbf{b}\}$ and $C$.
\end{enumerate}
\begin{proof}
\begin{enumerate}[label=(\roman*)]
\setlength{\itemsep}{0pt}
\item Put $m\defeq\inf\{\norm{\mathbf{x}}^2:\mathbf{x}\in C\}$. Choose $(\mathbf{x}_j)_{j\geq1}\subseteq C$ satisfying $m\leq\norm{\mathbf{x}_j}^2<m+\frac{1}{j}$ by the definition of infimum.\\
For $j,\ell\geq1$, the midpoint $\mathbf{w}_{j,\ell}\defeq\frac{1}{2}(\mathbf{x}_j+\mathbf{x}_\ell)$ lies in $C$ by convexity.\\
Hence $\norm{\mathbf{x}_j-\mathbf{x}_\ell}^2=2\norm{\mathbf{x}_j}^2+2\norm{\mathbf{x}_\ell}^2-4\norm{\mathbf{w}_{j,\ell}}^2\leq\frac{2}{j}+\frac{2}{\ell}$ by Theorem \ref{thm:parallelogram-identity-inner-product}, so $(\mathbf{x}_j)_{j\geq1}$ is Cauchy.\\
Write $\mathbf{x}_j=(x_{j,1},\ldots,x_{j,n})$. For $r\in\{1,\ldots,n\}$ and $q\geq1$, put $a_{q,r}\defeq\inf\{x_{j,r}:j\geq q\}$ and $b_{q,r}\defeq\sup\{x_{j,r}:j\geq q\}$.\\
These numbers exist by boundedness of each Cauchy coordinate sequence. The intervals $[a_{q,r},b_{q,r}]$ are nested, and $b_{q,r}-a_{q,r}\to0$: for $\eta>0$, choose $q_0$ with $\abs{x_{j,r}-x_{\ell,r}}<\eta$ whenever $j,\ell\geq q_0$ by the Cauchy property, so $0\leq b_{q,r}-a_{q,r}\leq\eta$ for $q\geq q_0$.\\
Put $N_r\defeq\sup\{a_{q,r}:q\geq1\}$. Then $a_{q,r}\leq N_r\leq b_{q,r}$ by nestedness, so $x_{j,r}\to N_r$.\\
For $\mathbf{N}_C\defeq(N_1,\ldots,N_n)$, $\norm{\mathbf{x}_j-\mathbf{N}_C}^2=\sum\limits_{r=1}^{n}\abs{x_{j,r}-N_r}^2\to0$. If $\mathbf{N}_C\notin C$, fix $\rho>0$ with $B_\rho(\mathbf{N}_C)\subseteq\R^n\setminus C$ by openness of $\R^n\setminus C$, contradicting $\mathbf{x}_j\to\mathbf{N}_C$; hence $\mathbf{N}_C\in C$. Also $\norm{\mathbf{x}_j}\to\norm{\mathbf{N}_C}$ by Proposition \ref{prop:reverse-triangle-inequality-inner-product}, so $\norm{\mathbf{N}_C}^2=m$.\\
If $\mathbf{u},\mathbf{v}\in C$ both have squared norm $m$, then $\mathbf{w}\defeq\frac{1}{2}(\mathbf{u}+\mathbf{v})\in C$, $m\leq\norm{\mathbf{w}}^2=m-\frac{1}{4}\norm{\mathbf{u}-\mathbf{v}}^2$ by Theorem \ref{thm:parallelogram-identity-inner-product}; thus $\mathbf{u}=\mathbf{v}$.
\item The singleton $\{-\mathbf{b}\}$ is convex directly and compact by Proposition \ref{prop:convex-hull-finite-set-combinations-compact}(i),(ii).\\
Since $C\neq\emptyset$, $D=\{-\mathbf{b}\}+C\neq\emptyset$; $D$ is closed, convex by Proposition \ref{prop:minkowski-sum-convex-compact-closed}(i),(ii). Fix its unique minimum-norm vector $\mathbf{N}$ by (i).\\
The condition $\mathbf{b}\notin C$ is equivalent to $\mathbf{0}\notin D$, so $\mathbf{N}\neq\mathbf{0}$. Fix $\mathbf{z}\in D$ and $t\in(0,1]$.\\
Then $\mathbf{N}+t(\mathbf{z}-\mathbf{N})\in D$ by convexity, and $0\leq\norm{\mathbf{N}+t(\mathbf{z}-\mathbf{N})}^2-\norm{\mathbf{N}}^2=2t\innerprod{\mathbf{N},\mathbf{z}-\mathbf{N}}+t^2\norm{\mathbf{z}-\mathbf{N}}^2$ by minimality.\\
If $a\defeq\innerprod{\mathbf{N},\mathbf{z}-\mathbf{N}}<0$, then $\mathbf{z}\neq\mathbf{N}$ and $2ta+t^2\norm{\mathbf{z}-\mathbf{N}}^2<0$ for $t\in\left(0,\min\left\{1,\frac{-2a}{\norm{\mathbf{z}-\mathbf{N}}^2}\right\}\right)$, a contradiction.\\
Thus $\innerprod{\mathbf{N},\mathbf{z}}\geq\norm{\mathbf{N}}^2$. Taking $\mathbf{z}=\mathbf{x}-\mathbf{b}$ proves $(\forall\,\mathbf{x}\in C)\,\innerprod{\mathbf{N},\mathbf{x}}\geq\norm{\mathbf{N}}^2+\innerprod{\mathbf{N},\mathbf{b}}>\innerprod{\mathbf{N},\mathbf{b}}$.
\item Put $\varepsilon\defeq\frac{1}{4}\norm{\mathbf{N}}^2>0$. By (ii), $(\forall\,\mathbf{x}\in C)\,\innerprod{\mathbf{N},\mathbf{b}}<\alpha-\varepsilon<\alpha+\varepsilon<\innerprod{\mathbf{N},\mathbf{x}}$.\\
Thus $\mathcal{H}(\mathbf{N},\alpha)$ strongly separates $\{\mathbf{b}\}$ and $C$ by the defining condition with $\sigma=1$.
\end{enumerate}
\hlt{Structural Commentary}: Midpoint estimates form a Cauchy minimizing sequence; translation converts its limit into separation.\\
\hlt{Key Insight}: Minimal norm forces a supporting linear inequality.
\end{proof}
\end{theorem}

Figure \ref{fig:convex-point-separating-hyperplane} illustrates Theorem \ref{thm:nearest-point-and-point-separation-convex-set}(ii),(iii): $\mathbf{p}=\mathbf{b}+\mathbf{N}$ is the nearest point in $C$, and $\mathbf{N}$ is normal to a separating hyperplane. The theorem, rather than the diagram, owns these claims.

\begin{illustrationfigure}[H]
  \centering
  % Source-supported schematic: Roman, Advanced Linear Algebra, Chapter 15.
% A point outside a closed convex set is strongly separated by a hyperplane.
\begin{tikzpicture}[
  x=0.88cm,
  y=0.86cm,
  >=Stealth,
  every node/.style={font=\small},
  convex/.style={draw=teal!65!black, fill=teal!8, line width=0.8pt},
  hyper/.style={red!70!black, line width=0.95pt},
  normal/.style={->, red!70!black, line width=0.9pt},
  guide/.style={gray!60, line width=0.5pt}
]
  \path[use as bounding box] (-0.35,-0.75) rectangle (5.85,2.85);

  \coordinate (b) at (0.55,0.25);
  \coordinate (p) at (2.35,0.85);

  \draw[convex]
    (p)
    .. controls (2.60,1.68) and (3.34,2.42) .. (4.48,2.18)
    .. controls (5.48,1.96) and (5.70,0.72) .. (4.90,0.02)
    .. controls (4.02,-0.72) and (2.76,-0.28) .. (p)
    -- cycle;
  \node[text=teal!60!black] at (4.42,1.55) {$C$};

  \draw[hyper] (0.7167,2.75) -- (1.7767,-0.43)
    node[pos=0.13, right=2pt] {$\mathcal{H}$};

  \node[circle, fill=black, inner sep=1.4pt, label=below:{$\mathbf{b}$}] at (b) {};
  \node[circle, fill=black, inner sep=1.2pt, label=above right:{$\mathbf{p}$}] at (p) {};

  \draw[normal] (b) -- (p)
    node[pos=0.52, below=5pt] {$\mathbf{N}$};

  \draw[guide] (1.57,0.59) -- (1.53,0.71) -- (1.41,0.67);
\end{tikzpicture}

  \caption{For $\mathbf{b}\notin C$, the labeled point $\mathbf{p}=\mathbf{b}+\mathbf{N}\in C$ is nearest to $\mathbf{b}$, the arrow $\mathbf{N}$ is the displacement from $\mathbf{b}$ to $\mathbf{p}$, and the labeled hyperplane lies strictly between $\mathbf{b}$ and the closed convex set $C$.}
  \label{fig:convex-point-separating-hyperplane}
\end{illustrationfigure}

\begin{theorem} \label{thm:compact-convex-set-separated-from-subspace} \hlt{Separating a Compact Convex Set from a Subspace}\\
Let $C\subseteq\R^n$ be nonempty, compact, and convex, and let $S\leq\R^n$ satisfy $C\cap S=\emptyset$. Then:
\begin{enumerate}[label=(\roman*)]
\setlength{\itemsep}{0pt}
\item $(\exists\,\mathbf{N}\in S^\perp\setminus\{\mathbf{0}\})\,\innerprod{\mathbf{N},C}\geq\norm{\mathbf{N}}^2>0$.
\item For a vector $\mathbf{N}$ as in (i), $\mathcal{H}(\mathbf{N},\frac{1}{2}\norm{\mathbf{N}}^2)$ strongly separates $S$ and $C$.
\end{enumerate}
\begin{proof}
\begin{enumerate}[label=(\roman*)]
\setlength{\itemsep}{0pt}
\item The subspace $S$ is convex. Put $Q\defeq P_{S^\perp}$. The map $Q$ is linear and $\ker(Q)=S$ by Proposition \ref{prop:orthogonal-projection-properties}(i),(vi) and Corollary \ref{cor:double-orthogonal-complement}.\\
Fix $\mathbf{x}\notin S$ and put $\mathbf{q}\defeq Q\mathbf{x}\neq\mathbf{0}$. If $\norm{\mathbf{y}-\mathbf{x}}<\frac{1}{2}\norm{\mathbf{q}}$, then $\norm{Q\mathbf{y}}\geq\norm{\mathbf{q}}-\norm{Q(\mathbf{y}-\mathbf{x})}\geq\norm{\mathbf{q}}-\norm{\mathbf{y}-\mathbf{x}}>0$ by Proposition \ref{prop:reverse-triangle-inequality-inner-product} and Proposition \ref{prop:orthogonal-projection-properties}(i),(viii). Thus $B_{\frac{1}{2}\norm{\mathbf{q}}}(\mathbf{x})\subseteq\R^n\setminus S$, so $S$ is closed.\\
Put $D\defeq C+S$. Since $\mathbf{0}\in S$, $C\subseteq D$ and $D\neq\emptyset$; and $D$ is closed and convex by Proposition \ref{prop:minkowski-sum-convex-compact-closed}(i),(ii).\\
If $\mathbf{0}\in D$, choose $(\mathbf{c},\mathbf{s})\in C\times S$ with $\mathbf{0}=\mathbf{c}+\mathbf{s}$. Then $\mathbf{c}=-\mathbf{s}\in C\cap S$, a contradiction.\\
Fix the unique minimum-norm vector $\mathbf{N}$ of $D$ by Theorem \ref{thm:nearest-point-and-point-separation-convex-set}(i).\\
Then $\mathbf{N}\neq\mathbf{0}$ and $(\forall\,\mathbf{c}\in C)(\forall\,\mathbf{s}\in S)\,\innerprod{\mathbf{N},\mathbf{c}+\mathbf{s}}\geq\norm{\mathbf{N}}^2$ by Theorem \ref{thm:nearest-point-and-point-separation-convex-set}(ii) with the closed convex set $D$ and the point $\mathbf{0}$.\\
Fix $\mathbf{c}_0\in C$ and suppose $\mathbf{s}\in S$ satisfies $a\defeq\innerprod{\mathbf{N},\mathbf{s}}\neq0$. Set $\lambda\defeq\frac{\norm{\mathbf{N}}^2-\innerprod{\mathbf{N},\mathbf{c}_0}-1}{a}$. Then $\lambda\mathbf{s}+\mathbf{c}_0\in D$ but $\innerprod{\mathbf{N},\lambda\mathbf{s}+\mathbf{c}_0}=\norm{\mathbf{N}}^2-1$, a contradiction. Hence $(\forall\,\mathbf{s}\in S)\,\innerprod{\mathbf{N},\mathbf{s}}=0$, so $\mathbf{N}\in S^\perp$. Taking $\mathbf{s}=\mathbf{0}$ proves $\innerprod{\mathbf{N},C}\geq\norm{\mathbf{N}}^2>0$.
\item Put $\alpha\defeq\frac{1}{2}\norm{\mathbf{N}}^2$ and $\varepsilon\defeq\frac{1}{4}\norm{\mathbf{N}}^2$. By (i), $(\forall\,\mathbf{s}\in S)(\forall\,\mathbf{c}\in C)\,\innerprod{\mathbf{N},\mathbf{s}}=0<\alpha-\varepsilon<\alpha+\varepsilon<\norm{\mathbf{N}}^2\leq\innerprod{\mathbf{N},\mathbf{c}}$.\\
Thus $\mathcal{H}(\mathbf{N},\alpha)$ strongly separates $S$ and $C$ by the defining condition with $\sigma=1$.
\end{enumerate}
\hlt{Structural Commentary}: Separation of $C+S$ constructs a normal; scaling along $S$ proves orthogonality.\\
\hlt{Key Insight}: Translation invariance along $S$ forces the normal into $S^\perp$.
\end{proof}
\end{theorem}
